% !TeX program = xelatex
\documentclass[fontsize=14pt,a4paper]{scrartcl}
\usepackage{fontspec}
\usepackage{polyglossia}
\setmainlanguage{russian}
\IfFontExistsTF{Times New Roman}{\setmainfont{Times New Roman}}{\setmainfont{FreeSerif}}
\IfFontExistsTF{Arial}{\setsansfont{Arial}}{\setsansfont{FreeSans}}
\IfFontExistsTF{Courier New}{\setmonofont{Courier New}\newfontfamily\cyrillicfonttt{Courier New}}{\setmonofont{FreeMono}\newfontfamily\cyrillicfonttt{FreeMono}}
\usepackage[left=30mm,right=15mm,top=20mm,bottom=20mm]{geometry}
\usepackage{amsmath,amssymb}
\usepackage{setspace,graphicx,float,enumitem}
\usepackage[hidelinks]{hyperref}
\onehalfspacing
\setlength{\parindent}{1.25cm}
\emergencystretch=3em
\setkomafont{disposition}{\rmfamily\bfseries}
\setkomafont{section}{\fontsize{14}{17}\selectfont}
\setkomafont{subsection}{\fontsize{14}{17}\selectfont}
\setlist{itemsep=0pt,topsep=3pt}
\newcommand{\screen}[2]{\begin{figure}[H]\centering\includegraphics[width=\textwidth,height=.72\textheight,keepaspectratio]{images/#1}\caption{#2}\end{figure}}
\begin{document}
\begin{titlepage}
\thispagestyle{empty}
\begin{center}
Московский авиационный институт\\
(национальный исследовательский институт)\\[1.5\baselineskip]
Институт «Компьютерные науки и прикладная математика»
\vspace{30mm}

\textbf{Практическая работа}\\
\textbf{«Неклассические логики»}\\
\textbf{VII семестр}\\[1.5\baselineskip]
Нечёткие величины на основе\\логнормального распределения
\end{center}
\vfill
\begin{flushright}
Выполнили студенты:\\
Сетраков Фёдор\\
Коваленко Денис\\
Группа: М8О-410Б-23\\[5mm]
Оценка: \rule{35mm}{0.2pt}\\[3mm]
Дата: \rule{35mm}{0.2pt}
\end{flushright}
\vfill
\begin{center}\textbf{Москва — 2026}\end{center}
\end{titlepage}
\tableofcontents
\clearpage
\section*{Задание}
\addcontentsline{toc}{section}{Задание}
Разработать на Go программный комплекс для представления нечётких величин заданного вида, выполнения арифметических операций и сравнений, генерации наборов из 100--300 объектов и кластерного анализа. Реализовать разбиение на заданное число кластеров и разбиение по критерию, графический интерфейс, сохранение и загрузку данных, журналирование и тестирование. В настоящей работе используется логнормальная модель.

Расчётные процедуры отделены от интерфейса и оформлены пакетами. В текущей структуре пакеты находятся в \texttt{internal}: они доступны внутри родительского дерева проекта. Для прямого импорта сторонним Go-модулем потребуется вынести публичный API за пределы \texttt{internal}.

\section{Цель работы}

Изучить представление нечётких величин, арифметику через \(\alpha\)-срезы и зависимость кластерного разбиения от уровня принадлежности. Разработать независимые библиотеки на Go и графическое приложение для экспериментов с наборами из 100–300 величин.

\section{Вариант}

В работе используется логарифмически нормальное распределение с параметрами \(\mu\) и \(\sigma\) > 0:

\[
f(x;\mu,\sigma)=\begin{cases}
\dfrac{1}{x\sigma\sqrt{2\pi}}\exp\left[-\dfrac{(\ln x-\mu)^2}{2\sigma^2}\right],&x>0,\\
0,&x\leq0.
\end{cases}
\]

Это означает $\ln X\sim N(\mu,\sigma^2)$. Параметры не являются средним и отклонением самого $X$. Для исходной случайной величины:
\[
E[X]=e^{\mu+\sigma^2/2},\qquad
\operatorname{Var}[X]=(e^{\sigma^2}-1)e^{2\mu+\sigma^2},
\]
а мода равна $m=e^{\mu-\sigma^2}$.


Источник формулы исходного распределения: \href{https://itl.nist.gov/div898/handbook/eda/section3/eda3669.htm}{NIST, Lognormal Distribution}.

\section{Переход к функции принадлежности}

Исходная формула задаёт плотность вероятности; она может превышать единицу и не задаёт однозначно способ перехода к нечёткой величине. В программе принято нормирование плотности по максимуму. CDF не используется как функция принадлежности: она монотонна и стремится к единице при x\(\to\)\(\infty\), что приводит к неограниченным срезам.

Пусть $c=\mu-\sigma^2$. Деление плотности на её максимум даёт:

\[
g(x)=f(x)/f(e^c)=\exp\left[-\frac{(\ln x-c)^2}{2\sigma^2}\right],\qquad x>0.
\]

g имеет высоту 1, единственную моду exp(c) и вложенные интервальные \(\alpha\)-срезы. Однако её носитель бесконечен. Для численного расчёта выбран уровень \(\varepsilon\)=\(10^{-4}\) и непрерывная компактная модель:

\[
\widetilde\mu(x)=\max\left(0,\frac{g(x)-\varepsilon}{1-\varepsilon}\right).
\]

Таким образом, носитель — замкнутый интервал, на границах принадлежность равна нулю, высота остаётся единицей. Это допущение реализации. Для всех x абсолютное расхождение g и \(\widetilde\mu\) не превышает \(\varepsilon\); это \textbf{не} означает такой же относительной точности границ, моментов или вероятностей хвоста.

Методы Density и CDF возвращают значения исходного распределения без отсечения. Membership, Support и AlphaCut работают с компактной нечёткой моделью.

\section{Аналитические \(\alpha\)-срезы}

Обозначим $\beta=\varepsilon+(1-\varepsilon)\alpha$. Решение неравенства $\widetilde\mu(x)\geq\alpha$ при $0<\alpha\leq1$:

\[
A_\alpha=\left[\exp(c-\sigma\sqrt{-2\ln\beta}),\exp(c+\sigma\sqrt{-2\ln\beta})\right].
\]

При \(\alpha\)=0 принимается замыкание носителя с той же формулой для \(\beta\)=\(\varepsilon\). При \(\alpha\)=1 обе границы равны exp(c). С ростом \(\alpha\) левая граница возрастает, правая убывает. В отличие от равномерного варианта исходного проекта изменяются обе границы.

\section{Представление данных}

Интерфейс \texttt{fuzzy.Number} предоставляет функцию принадлежности, носитель, \(\alpha\)-срез, высоту и имя типа. \texttt{fuzzy.LogNormal} хранит параметры \(\mu\) и \(\sigma\). Конструктор проверяет конечность параметров, \(\sigma\)>0, допустимость диапазона экспоненты и различимость границ в машинной арифметике.

\texttt{fuzzy.Cuts} хранит сетку \(\alpha\) и вложенные интервалы. Между узлами границы интерполируются линейно. Этот тип позволяет хранить результат, даже когда он не принадлежит логнормальному семейству. \texttt{fuzzy.Collection} реализует набор величин: добавление, удаление, копирование, фильтрацию и общие границы.

\section{Арифметические операции}

По принципу обобщения Заде операции выполняются над \(\alpha\)-срезами. Для интервалов $A=[l_1,r_1]$, $B=[l_2,r_2]$:

\begin{itemize}
\item $A+B=[l_1+l_2,r_1+r_2]$.
\item $A-B=[l_1-r_2,r_1-l_2]$.
\item A\(\times\)B определяется минимумом и максимумом четырёх произведений концов.
\item $A/B$ определяется интервальным умножением на $[1/r_2,1/l_2]$, если $0\notin B$.
\end{itemize}

Для каждого уровня вычисляется результирующий интервал; результат сохраняется как Cuts. По умолчанию используется 200 шагов. Чем сильнее нелинейность границ, тем существеннее влияние разрешения. Это арифметика нечётких величин, а не вероятностная свёртка.

Разность положительных исходных величин может содержать отрицательные значения. У допустимой LogNormal носитель строго положителен, поэтому прямое деление двух исходных величин разрешено. Деление на произвольный результат с нулём в носителе отклоняется.

\section{Сравнение и характеристики}

Возможность Pos(A\(\geq\)B) — максимальный уровень \(\alpha\), на котором правая граница A\(\alpha\) не меньше левой границы B\(\alpha\). Значение находится двоичным поиском. Необходимость для нормальных величин реализована как Nec(A\(\geq\)B)=1-Pos(B>A). Здесь используется строгое обратное отношение: Pos(B>A) — супремум уровней \(\alpha\), на которых правая граница B\(\alpha\) строго больше левой границы A\(\alpha\). Супремум может не достигаться; для одинаковых чётких величин Pos(B>A)=0 и Nec(A\(\geq\)B)=1.

Библиотека также предоставляет стратегии ранжирования по центру тяжести, середине среза и разности возможностей. Ранжирование по возможности может давать равенство сильно пересекающихся величин.

Площадь: $S=\int_0^1(R_\alpha-L_\alpha)\,d\alpha$. Центр тяжести:

\[
C=\frac{\int_0^1(R_\alpha^2-L_\alpha^2)/2\,d\alpha}{\int_0^1(R_\alpha-L_\alpha)\,d\alpha}.
\]

Возможностные характеристики Карлссона–Фуллера:

\[
E_f=\int_0^1\alpha(L_\alpha+R_\alpha)d\alpha,\qquad
V_f=\frac12\int_0^1\alpha(R_\alpha-L_\alpha)^2d\alpha.
\]

Интегралы вычисляются методом трапеций. Эти величины не равны вероятностным моментам X. Интерфейс помечает их отдельно, чтобы избежать смешения понятий.

Библиотека расстояний включает интегральные L1/L2 по границам срезов, максимум расстояния Хаусдорфа по \(\alpha\)-сетке, расстояние центроидов, интеграл |\(\mu\)A-\(\mu\)B| по x и метрику одного \(\alpha\)-среза.

\section{Генерация}

Параметры \(\mu\) и \(\sigma\) выбираются независимо равномерно в заданных диапазонах. Режим «облака» сначала выбирает K центров в пространстве (\(\mu\),\(\sigma\)), затем добавляет равномерный шум в пределах доли диапазона. Параметры ограничиваются заданными границами; \(\sigma\) остаётся положительным. Объекты распределяются между центрами по очереди. Облака не гарантируют точного восстановления исходных групп кластеризатором.

Seed обеспечивает воспроизводимость. Интерфейс ограничивает генерацию 100–300 величинами; библиотека допускает 1–300 для небольших примеров и тестов.

\section{Кластерный анализ}

В интерфейсе используется расстояние одного среза:

\[
d_\alpha(A,B)=\frac{|L_A(\alpha)-L_B(\alpha)|+|R_A(\alpha)-R_B(\alpha)|}{2}.
\]

На его основе строится симметричная матрица попарных расстояний. Алгоритмы кластеризации работают с матрицей и не зависят от вида нечёткой величины:

\begin{enumerate}
\item \textbf{PAM}. Жадно выбирает k начальных медоидов, затем заменяет медоиды другими объектами, пока суммарное расстояние до ближайших медоидов уменьшается.
\item \textbf{Иерархическая кластеризация}. Начинает с отдельных объектов, объединяет ближайшие группы со средней связью; останавливается, когда минимальное расстояние достигает порога. Количество групп определяется порогом.
\end{enumerate}

Медоид — реальный объект набора, минимизирующий сумму расстояний до членов группы. При \(\alpha\)=1 расстояния определяются только модами exp(\(\mu\)-\(\sigma^2\)); объекты с одинаковой модой и разными \(\sigma\) становятся неразличимы выбранной метрикой. При \(\alpha\)<1 учитывается также ширина.

Память для матрицы O(\(N^2\)); используемые прямые реализации алгоритмов ориентированы на учебные наборы до 300 объектов.

\section{Архитектура}

Все Go-пакеты, кроме \texttt{cmd}, расположены на одном уровне внутри \texttt{internal}. Пакет \texttt{distribution} содержит вероятностную модель и её моменты; \texttt{fuzzy} — нечёткую модель, операции и метрики. \texttt{generator} создаёт наборы; \texttt{clustering} работает только с матрицами расстояний. \texttt{storage} отвечает за CSV, \texttt{logger} — за журнал.

\texttt{experiment.Service} управляет состоянием, выполнением сценариев и ограниченным кешем результатов. Он не зависит от HTTP и геометрии графиков. \texttt{web} разбирает запросы, вызывает сервис и строит HTML/SVG. \texttt{cmd} создаёт зависимости и запускает сервер. GoDoc написан на английском.

Статические ресурсы встроены через embed. Генерация открывает предварительный просмотр без кластеризации; расчёт запускается явно. Переключение вкладок графика не выполняет расчётов. Сервис возвращает копии снимков и защищает общее состояние mutex-ом. Запись нового набора происходит только после успешного журналирования. Для импортированных CSV seed неизвестен.

\section{Руководство пользователя}

\begin{enumerate}
\item Установить Go \(\geq\)1.24.7, перейти в каталог проекта, выполнить \texttt{go run ./cmd}.
\item Открыть \texttt{http://localhost:8080}.
\item Задать диапазоны \(\mu\) и \(\sigma\), количество 100–300, режим и seed. Нажать «Создать и посмотреть набор». Изучить плотность, функцию принадлежности и исходные данные.
\item Выбрать PAM и k либо иерархический метод и положительный порог.
\item Задать \(\alpha\) и нажать «Запустить кластеризацию». Изменение параметров само по себе не запускает вычисления.
\item Изучить срезы, карту границ и медоиды. Таблица «Данные и вероятностные характеристики» содержит численные данные графиков.
\item Для CSV открыть «Сохранение и загрузка». Подготовленный CSV нужно поместить в \texttt{data/input} и указать его имя.
\item На странице «Операции» задать параметры A и B, выбрать операцию и нажать «Рассчитать».
\item Переключатель в шапке меняет тему и сохраняет выбор в браузере.
\end{enumerate}

CSV содержит заголовок и строки параметров следующего вида:
\begin{samepage}
\begin{verbatim}
kind,mu,sigma
lognormal,0,0.3
lognormal,0.7,0.4
\end{verbatim}
\end{samepage}
 Сохраняются исходные параметры без округления. Загрузчик проверяет формат, параметры, непустоту и ограничение размера. Журнал \texttt{data/log.json} содержит JSON Lines с временными метками, параметрами генерации и операций, размерами кластеров и ошибками вычислений.

\section{Проверка}

Команды: \texttt{go test ./...}, \texttt{go vet ./...}. Тесты проверяют исходную плотность и CDF, нормирование, совпадение границ аналитических срезов с Membership, вложенность, некорректные параметры, интервальную арифметику, сравнение, генерацию, оба метода кластеризации, CSV, журнал и HTTP-интерфейс.

Проверки интерфейса и скриншоты приведены в \texttt{docs/qa}. Они дополняют модульные тесты проверками htmx, переключения темы, полей методов, валидации и адаптивности.

\section{Ограничения и выводы}

Отсечение и нормирование являются выбранным способом получения нечёткой модели из заданной плотности. Их необходимо указывать при защите работы. Они не меняют исходные методы Density и CDF. Точность характеристик и арифметики зависит от численной сетки; \(\varepsilon\) не является универсальной оценкой их ошибки.

Программа предоставляет воспроизводимые эксперименты с логнормальными нечёткими величинами и показывает, как изменение \(\alpha\) меняет представление объектов и разбиение. Благодаря интерфейсу Number арифметика, метрики и кластеризация используются независимо от веб-приложения.

\section{Демонстрация работы программы}
Снимки выполнены 15 ноября 2026 года в текущей реализации приложения (коммит \texttt{b3dbee8}). Используется светлая тема интерфейса. На всех снимках эксперимента рассматривается один набор: 200 объектов, четыре облака, шум 0,03, seed 1, диапазоны $\mu\in[-1;2]$ и $\sigma\in[0{,}15;0{,}6]$.

\subsection{Предварительный просмотр}
После создания набора программа показывает исходные плотности и функции принадлежности первых трёх величин; полный набор доступен в таблице. Это ограничение визуализации, а не размера набора. Кластеризация использует все 200 объектов.
\screen{preview.jpg}{Предварительный просмотр: плотность вероятности и функция принадлежности}

\subsection{Кластеризация и представление интервалов}
При $\alpha=0$ используются замыкания рабочих носителей. Для PAM выбрано $k=4$. В демонстрационном запуске получены четыре группы по 50 объектов. Параметры медоидов, округлённые в интерфейсе:
\begin{enumerate}
\item $\mu=0{,}2867$, $\sigma=0{,}4594$; носитель примерно $[0{,}15;7{,}75]$.
\item $\mu=0{,}8065$, $\sigma=0{,}5759$; носитель примерно $[0{,}136;19]$.
\item $\mu=-0{,}8207$, $\sigma=0{,}2206$; носитель примерно $[0{,}163;1{,}08]$.
\item $\mu=0{,}9663$, $\sigma=0{,}3521$; носитель примерно $[0{,}512;10{,}5]$.
\end{enumerate}
Это наблюдение одного воспроизводимого запуска, а не гарантия восстановления облаков на любых параметрах. Число объектов в группах само по себе не является оценкой качества кластеризации.

На графике срезов одна горизонтальная полоска соответствует одному объекту. Её концы равны $L_\alpha$ и $R_\alpha$, а цвет обозначает найденный кластер. Вертикальное положение служит для размещения объектов и не является степенью принадлежности.
\screen{cuts.jpg}{PAM: интервалы 200 объектов при $\alpha=0$, четыре кластера}

На плоскости границ тот же интервал изображается точкой $(L_\alpha,R_\alpha)$. Например, интервал $[3;7]$ представлен точкой $(3,7)$. Медоиды выделены обводкой. Масштабы осей могут различаться, поэтому визуальное евклидово расстояние на рисунке не следует отождествлять с используемой метрикой.
\screen{plane.jpg}{Плоскость границ того же набора; медоиды выделены обводкой}

\subsection{Арифметика и сравнение}
В демонстрации заданы $A=\operatorname{LN}(0;0{,}3)$ и $B=\operatorname{LN}(0{,}7;0{,}4)$, выбрано сложение. Полученный рабочий носитель суммы приблизительно равен $[0{,}56;12{,}86]$. Функции A, B и результата отображаются различными типами линий. В таблице представлены носитель, центр тяжести, возможностные среднее и дисперсия.

Для исходных операндов интерфейс показывает $\operatorname{Pos}(A\ge B)\approx0{,}667$ и $\operatorname{Nec}(A\ge B)=0$. Следовательно, отношение допускается с ненулевой степенью возможности, но не является необходимым. Это степени нечёткого сравнения, а не вероятности событий.
\screen{operations.jpg}{Сложение двух логнормальных нечётких величин и характеристики результата}
\clearpage

\section{План исследования влияния уровня принадлежности}
Для сравнения результатов необходимо сохранять один исходный набор и настройки метода. Выполняется расчёт для уровней $\alpha=0$, $0{,}5$, $0{,}9$ и $1$. Сравниваются размеры групп, параметры медоидов и принадлежность объектов к группам. При $\alpha=1$ остаются только моды, поэтому информация о ширине исходных кривых не участвует в выбранном расстоянии. Данный план описывает следующий эксперимент; численные результаты для этих дополнительных уровней здесь не заявляются.

\section*{Заключение}
\addcontentsline{toc}{section}{Заключение}
Реализован комплекс для генерации и исследования логнормальных нечётких величин, арифметики по срезам, сравнения и кластерного анализа. Вероятностная и нечёткая модели разделены явно. Интерфейс позволяет прослеживать переход от параметров к функциям, интервалам и кластерам. Демонстрационный запуск подтверждает работу полного сценария на наборе из 200 объектов. При интерпретации результатов необходимо учитывать выбранный уровень среза, отсечение хвостов и численное разрешение.
\end{document}
